\section{Introducción}

	El análisis de potencia en sistemas eléctricos trifásicos permite determinar cómo se distribuye la energía entre la parte útil, la parte asociada a campos eléctricos o magnéticos y la demanda total que debe suministrar la fuente. En esta práctica se estudió una alimentación trifásica senoidal de frecuencia \(60\unit{Hz}\), con fases desfasadas \(120^{\circ}\), para observar el comportamiento de distintas cargas y verificar experimentalmente la influencia del factor de potencia en el consumo de corriente.
	
	El desarrollo contempla tres partes principales: primero, la revisión de las fuentes de alimentación de sistemas trifásicos mediante transformadores; segundo, el estudio de la potencia eléctrica en régimen sinusoidal para cargas resistivas, inductivas y capacitivas; y tercero, el análisis del factor de potencia, su medición y su corrección mediante un banco de capacitores. Con los datos de laboratorio se comparan las condiciones de operación del motor de forma aislada y posteriormente con compensación capacitiva.